% GENERATED FILE. Edit canonical lessons, then run npm run generate:books.
% canonical-source-hash: fnv1a64:415c38f77d18de94
% canonical-lessons: GU-C131-nicu, GU-C131-paholu, GU-C131-patlu, GU-C131-bhinu, GU-C131-suku

\chapter{Low, Wide, Thin, Wet, Dry}
\label{ch:gu-low-wide-thin-wet-dry}
\hlchaptermodality{131}

\begin{chapteropening}
\textbf{By the end of this chapter:} \emph{I can say low, wide, thin, wet and dry.}

Complete the last lesson of chapter 131: I can say low, wide, thin, wet and dry.
\end{chapteropening}

\section[\texorpdfstring{\gu{નીચું} (nīcũ)}{નીચું (nīcũ)}]{\gu{નીચું} (nīcũ) — low}

\label{lesson:GU-C131-nicu}



\begin{quote}
Before the new one: say the Gujarati for a salary, then the Gujarati for a tailor.
\end{quote}

\subsection*{You'll want to know: \gu{નીચું}}
\textbf{\gu{નીચું}} (\emph{nīcũ}) — \textquotedblleft{}low\textquotedblright{}.

\begin{grammarlens}[title={What you've built}]
One more word for this chapter.
\end{grammarlens}

\subsection*{Your turn}
\begin{itemize}
\raggedright
  \item \emph{Say it:} \emph{nīcũ}
  \item \emph{Say it:} \emph{nīcũ}, once more
  \item \emph{Say it:} say \emph{nīcũ}
  \item \emph{Recall:} say the Gujarati for a salary, then the Gujarati for a tailor, then say \emph{nīcũ} again
\end{itemize}

\begin{tcolorbox}[breakable,colback=teal!4,colframe=teal!35!black,title={Before you move on}]
What does \gu{નીચું} mean? (\textquotedblleft{}Low\textquotedblright{}.) Say it once more.
\end{tcolorbox}
\section[\texorpdfstring{\gu{પહોળું} (pahoḷũ)}{પહોળું (pahoḷũ)}]{\gu{પહોળું} (pahoḷũ) — wide}

\label{lesson:GU-C131-paholu}



\begin{quote}
Before the new one: say the Gujarati for a tailor, then the Gujarati for low.
\end{quote}

\subsection*{You'll want to know: \gu{પહોળું}}
\textbf{\gu{પહોળું}} (\emph{pahoḷũ}) — \textquotedblleft{}wide\textquotedblright{}.

\begin{grammarlens}[title={What you've built}]
One more word for this chapter.
\end{grammarlens}

\subsection*{Your turn}
\begin{itemize}
\raggedright
  \item \emph{Say it:} \emph{pahoḷũ}
  \item \emph{Say it:} \emph{pahoḷũ}, once more
  \item \emph{Say it:} say \emph{pahoḷũ}
  \item \emph{Recall:} say the Gujarati for a tailor, then the Gujarati for low, then say \emph{pahoḷũ} again
\end{itemize}

\begin{tcolorbox}[breakable,colback=teal!4,colframe=teal!35!black,title={Before you move on}]
What does \gu{પહોળું} mean? (\textquotedblleft{}Wide\textquotedblright{}.) Say it once more.
\end{tcolorbox}
\section[\texorpdfstring{\gu{પાતળું} (pātḷũ)}{પાતળું (pātḷũ)}]{\gu{પાતળું} (pātḷũ) — thin}

\label{lesson:GU-C131-patlu}



\begin{quote}
Before the new one: say the Gujarati for low, then the Gujarati for wide.
\end{quote}

\subsection*{You'll want to know: \gu{પાતળું}}
\textbf{\gu{પાતળું}} (\emph{pātḷũ}) — \textquotedblleft{}thin\textquotedblright{}.

\begin{grammarlens}[title={What you've built}]
One more word for this chapter.
\end{grammarlens}

\subsection*{Your turn}
\begin{itemize}
\raggedright
  \item \emph{Say it:} \emph{pātḷũ}
  \item \emph{Say it:} \emph{pātḷũ}, once more
  \item \emph{Say it:} say \emph{pātḷũ}
  \item \emph{Recall:} say the Gujarati for low, then the Gujarati for wide, then say \emph{pātḷũ} again
\end{itemize}

\begin{tcolorbox}[breakable,colback=teal!4,colframe=teal!35!black,title={Before you move on}]
What does \gu{પાતળું} mean? (\textquotedblleft{}Thin\textquotedblright{}.) Say it once more.
\end{tcolorbox}
\section[\texorpdfstring{\gu{ભીનું} (bhīnũ)}{ભીનું (bhīnũ)}]{\gu{ભીનું} (bhīnũ) — wet}

\label{lesson:GU-C131-bhinu}



\begin{quote}
Before the new one: say the Gujarati for wide, then the Gujarati for thin.
\end{quote}

\subsection*{You'll want to know: \gu{ભીનું}}
\textbf{\gu{ભીનું}} (\emph{bhīnũ}) — \textquotedblleft{}wet\textquotedblright{}.

\begin{grammarlens}[title={What you've built}]
One more word for this chapter.
\end{grammarlens}

\subsection*{Your turn}
\begin{itemize}
\raggedright
  \item \emph{Say it:} \emph{bhīnũ}
  \item \emph{Say it:} \emph{bhīnũ}, once more
  \item \emph{Say it:} say \emph{bhīnũ}
  \item \emph{Recall:} say the Gujarati for wide, then the Gujarati for thin, then say \emph{bhīnũ} again
\end{itemize}

\begin{tcolorbox}[breakable,colback=teal!4,colframe=teal!35!black,title={Before you move on}]
What does \gu{ભીનું} mean? (\textquotedblleft{}Wet\textquotedblright{}.) Say it once more.
\end{tcolorbox}
\section[\texorpdfstring{\gu{સૂકું} (sūkũ)}{સૂકું (sūkũ)}]{\gu{સૂકું} (sūkũ) — dry}

\label{lesson:GU-C131-suku}



\begin{quote}
Before the new one: say the Gujarati for thin, then the Gujarati for wet.
\end{quote}

\subsection*{You'll want to know: \gu{સૂકું}}
\textbf{\gu{સૂકું}} (\emph{sūkũ}) — \textquotedblleft{}dry\textquotedblright{}.

\begin{grammarlens}[title={What you've built}]
One more word for this chapter.
\end{grammarlens}

\subsection*{Your turn}
\begin{itemize}
\raggedright
  \item \emph{Say it:} \emph{sūkũ}
  \item \emph{Say it:} \emph{sūkũ}, once more
  \item \emph{Say it:} say \emph{sūkũ}
  \item \emph{Recall:} say the Gujarati for thin, then the Gujarati for wet, then say \emph{sūkũ} again
\end{itemize}

\begin{tcolorbox}[breakable,colback=teal!4,colframe=teal!35!black,title={Before you move on}]
What does \gu{સૂકું} mean? (\textquotedblleft{}Dry\textquotedblright{}.) Say it once more.
\end{tcolorbox}
